\newpage
\chapter{申請審査}

\section{審査に臨む姿勢}

申請審査は活動を選別する関門ではない。
この申請を認めてよいかではなく、この申請を成立させるために何が必要かを考える。

認可しないという判断は、確認と補正を尽くしてなお成立しない場合に限られる。

活動内容の優劣は判断しない。
題材が地味に見えても、対象とする音楽ゲームが少数派でも、審査で考慮しない。

\section{学生申請}

学生資格は申請をもって成立するため（学園規則第４章第３条）、認可の対象ではなく、審査を行わない。

確認するのは、学園が活動の基盤として利用する外部サービスの利用条件を満たしているかのみである（学園規則第４章第２条第２項）。
これは学園独自の制限ではなく外部条件の反映であり、ここに裁量はない。

\section{講師登録申請の確認事項}

担当しようとする分野について、他の学生に伝えられるだけの知識と経験があるかを実質的に確認する。
形式確認にとどめないのは、講義内容の信頼性が学園全体の信頼に関わるためである。

ただし高い水準を求めるという意味ではない。
見るのは、語ろうとしていることに裏づけがあるかどうかであり、競技的な上位者であることや経験年数の長さは要件ではない。

\section{担当分野が狭い場合の取扱い}

担当分野は広く宣言しておくほうが、その後の活動の自由度が高くなる。
申請内容が狭すぎるときは、より広く書くよう提案する。

これは審査上の指摘ではなく、申請者の利益になる助言として伝える。

\section{同意事項の確認}

学園規則第６章第２条第１項の同意が申請に含まれているかを確認する。

\begin{enumerate}[label=(\arabic*)　]
  \item 学内活動の目的に必要な範囲での講義資料の無償使用及び配信
  \item 講師に関するガイドラインの遵守
\end{enumerate}

同意がなければ講師となる要件を欠くため、申請は成立しない。

\section{講義開講申請}

講義開講申請は原則としてすべて認可する。
講義の内容、構成及び進め方は講師の裁量に属する（学園規則第６章第３条）。

認可しないのは次の場合に限る。

\begin{enumerate}[label=(\arabic*)　]
  \item 記入漏れその他の形式不備があるとき
  \item 明らかに虚偽の情報を伝える意図がうかがえるとき
\end{enumerate}

第１号は差し戻して補正を求めるものであり、不認可ではない。
第２号は、学園規則第７章第６条第１項第１号が講義の中止又は変更の事由として挙げるものと同種のものである。

内容が独特であること、範囲が狭いこと、既存の講義と重複することは、認可しない理由にならない。

\section{区分が内容と合わない場合の取扱い}

申請ガイドラインに定める区分が講義内容と合っていないと思われるときは、選び直しを提案する。

これは形式不備の差戻しではなく助言として扱い、申請者が当初の区分を維持したいと述べたときはそのまま認可する。
区分は学生が講義を探すための目印であり、内容の当否を判断するためのものではない。

\section{講義回数}

講義回数は講師が予定する回数であり、審査で増減を求めない。

学園規則第７章第５条は、予定した回数を正規活動期間内に実施できないときに補講期間で補うことを認めている。
したがって回数は達成義務ではなく、多すぎるか少なすぎるかを審査で判断しない。

講義の構成は講師の裁量に属する（学園規則第６章第３条）。

\section{担当分野と音楽ゲームとの接続}

担当分野そのものが音楽ゲームでなくても、音楽ゲームを目的とした講義であれば認可する。
学園憲章第２条は、活動の形式や対象を問わず、音楽ゲームに向き合う意志を持つ者の探究活動を等しく支援することを定めている。

判断材料は担当分野の名称ではなく、担当分野の選定理由である。
そこに音楽ゲームとの接続が書かれていれば足りる。
接続が読み取れないときは、不認可とする前に選定理由の補足を求める。

\section{実績の記載がない場合の取扱い}

実績は任意の記載事項であり、記載がないことを理由に不認可としない。

確認するのは担当分野についての裏づけであり、これは選定理由その他の記載からも読み取れる。
実績の記載量で判断しない。

\section{不認可の判断}

講師登録申請及び講義開講申請の認可は、教務主事の全員一致による（学園規則第６章第２条第２項及び第７章第２条第２項）。
単独で不認可を決めない。

疑義が生じたときは、まず申請者に確認する。
多くは説明や補正で解消するが、解消しないときは他の教務主事に共有し、取扱いは学園長が取りまとめる。

不認可と決したときは理由を申請者に伝える（学園規則第５章第６条第２項）。
理由を述べられない判断は行わない。

不認可について申請者から改めて申出があったときは、当該判断を行った者以外の教務主事が確認する。

\section{記録}

審査の経緯を記録に残す。
特に、確認や補正を求めたとき、不認可としたときは、その根拠を残す。

教務主事の職務は交代を超えて継承される（学園規則第５章第７条）。
記録がなければ、次に同種の申請が来たときに別の判断がなされ、申請者から見て一貫性のない運用になる。
